\documentclass[12pt,a4paper]{article}

% --- Преамбула ---
\usepackage[utf8]{inputenc}
\usepackage[T1]{fontenc}
\usepackage[russian]{babel}
\usepackage{amsmath,amssymb}
\usepackage{hyperref}
\usepackage[margin=2.5cm]{geometry}
\usepackage{enumitem}
\usepackage{setspace}
\usepackage{booktabs}
\onehalfspacing

\title{\textbf{Дополнение к Субквантовой теории}\\
\large Принцип масштабной проницаемости}
\author{}
\date{}

\begin{document}
\maketitle

\section{Определения}

\begin{quote}
\textbf{Определение 0.2.9 (Площадь участка спектра).} Для любой моды $M_{A \to B}$ и любого интервала $[k_1, k_2]$ определена \textbf{площадь участка спектра}:
\begin{equation}
    S_{A \to B}(k_1, k_2) = \sum_{k=k_1}^{k_2} |F_{A \to B}(k)|
\end{equation}
Это произведение средней амплитуды спектра на длину участка.
\end{quote}

\begin{quote}
\textbf{Определение 0.2.10 (Принцип масштабной проницаемости).} Если для двух мод $M_{A \to B}$ и $M_{C \to D}$ и некоторых участков $[k_1, k_2]$ и $[k'_1, k'_2]$ выполняется:
\begin{equation}
    S_{A \to B}(k_1, k_2) \gg S_{C \to D}(k'_1, k'_2),
\end{equation}
то корреляция между этими участками пренебрежимо мала:
\begin{equation}
    \sum_{k=k_1}^{k_2} F_{A \to B}(k) \cdot F_{C \to D}(k'_2 - k) \approx 0.
\end{equation}
\end{quote}

\section{Физический смысл}

Частица (или поле) с большей площадью участка спектра «не замечает» частицу с меньшей площадью. Большой участок проскакивает сквозь малый, как длинная волна проходит сквозь мелкое отверстие без рассеяния.

\textbf{Следствия:}
\begin{itemize}
    \item \textbf{Нейтрино} (очень малая площадь спектра) проходят сквозь обычную материю (большая площадь), почти не взаимодействуя.
    \item \textbf{Гравитация} (огромная длина, малая амплитуда) проходит сквозь всё, но её эффект накапливается на больших масштабах.
    \item \textbf{Сильное взаимодействие} (короткий, но высокоамплитудный пик) не проникает далеко, но доминирует на малых расстояниях.
    \item \textbf{Электромагнитное взаимодействие} (средняя площадь) занимает промежуточное положение.
\end{itemize}

\section{Связь с проницаемостью плато}

Плато — участок спектра с малой амплитудой, но потенциально большой длиной. Его площадь может быть как больше, так и меньше площади пика другой частицы. Это объясняет, почему:

\begin{itemize}
    \item Плато электрона (инертная масса) не взаимодействует с быстрыми частицами (их пики имеют меньшую площадь и проскакивают).
    \item Плато протона (кулоновский барьер) проницаемо для нейтронов (их спектр имеет большую площадь за счёт цветовых нулей) и непроницаемо для протонов (их спектры имеют сравнимую площадь, и корреляция велика).
\end{itemize}

\section{Формулировка в терминах собственных мод}

В спектральном представлении (Часть~2 Субквантовой теории) принцип масштабной проницаемости эквивалентен утверждению:

\begin{quote}
Если $\lambda_n \ll \lambda_m$, то $M_{nm}^{AB}(d) \approx 0$ для любых $d$.
\end{quote}

То есть собственные моды с сильно различающимися собственными значениями (масштабами) не коррелируют.

\section{Экспериментальная проверка}

Принцип может быть проверен в экспериментах по рассеянию:

\begin{itemize}
    \item Пучок частиц с известной площадью спектра направляется на мишень.
    \item Выход взаимодействия измеряется как функция отношения площадей пучка и мишени.
    \item \textbf{Предсказание:} выход резко падает, когда площадь пучка значительно больше площади мишени (или наоборот).
\end{itemize}

\end{document}